\section{刑法论 - 刑法的解释}
\begin{enumerate}
  \item 研究方向。立法论研究 + 解释论/刑法解释学/刑法教义学（实务）。
  \item 渊源。刑法典、刑法修正案（11 个）、单行法 1 部。
  \item 解释方法。
        \begin{enumerate}
          \item \zd{解释技巧}（生产部门）。生产结论。
                \begin{enumerate}
                  \item 平义解释。字面解释。不扩大不缩小。\al{公私财物解释为他人的财物}。
                  \item 扩大/扩张解释。解释后的含义比字面意义要大。\al{运钞车解释成银行}。\zd{把碗洗了}。
                  \item 缩小解释。和扩大解释相反。\al{在公共交通工具上抢劫、在摩的上抢劫}。
                  \item 类推解释。
                        \begin{enumerate}
                          \item \zd{禁止类推解释}。例外。有利于被告人。\al{强奸罪的妇女解释成包括男性}。
                          \item \zd{vs 扩大解释}。程度之分，无本质之分。\zd{扯得还行，扯得太远了}。
                        \end{enumerate}
                  \item 反义解释。从反面去解释。\al{年满 14 对抢劫罪负责，不满 14 无责}。
                \end{enumerate}
          \item 解释理由（质检部门）。论证结论是否合理。
                \begin{enumerate}
                  \item 文理/文义解释。字词、概念、术语的文字字义的解释。\al{丈夫强奸妻子能构成强奸罪}。
                  \item \zd{体系解释}。不同条文同一用语可以做不同解释。
                        \begin{enumerate}
                          \item 一词多义。\al{传播性病，传播淫秽物品；收买被拐卖妇女儿童，收买警察}。
                          \item 多词一义。\zd{恐吓 = 威胁 = 胁迫 = 敲诈} \yjt 以恶害相通告，使对方产生恐惧心理。
                        \end{enumerate}
                  \item 同类解释。针对兜底条款。
                  \item 当然解释。针对有罪无罪，是一种逻辑推理\zd{比较}，比较的两者\zd{性质相同程度不同}，\zd{要具备可比性}。
                        \begin{enumerate}
                          \item 举重以明轻 \yjt 论证无罪。举轻以明重 \yjt 论证有罪。
                          \item 抢劫强奸（不能）、强奸猥亵（能）。\al{和马云比有钱，和吴彦祖比帅}。
                        \end{enumerate}
                  \item 目的解释。根据罪\zd{要保护的法益}去解释。非法侵入住宅罪，法益许可权 \yjt 住宅的安宁权。
                \end{enumerate}
          \item 解释技巧和解释理由关系。\zd{相辅相成，相互制约}。
        \end{enumerate}
\end{enumerate}
